\begin{thebibliography}{62}

\bibitem[1]{AZGLOW}
Zeyuan Allen-Zhu, Ankit Garg, Yuanzhi Li, Rafael Oliveira, and Avi Wigderson. Operator scaling via geodesically convex optimization, invariant theory and polynomial identity testing. In STOC’18—Proceedings of the 50th Annual ACM SIGACT Symposium on Theory of Computing, pages 172–181. ACM, New York, 2018. doi:10.1145/3188745.3188942.

\bibitem[2]{audin2012torus}
Michele Audin. Torus actions on symplectic manifolds, volume 93. Birkhäuser, 2012.

\bibitem[3]{BDS}
Eric Bach, James R. Driscoll, and Jeffrey O. Shallit. Factor refinement. In David S. Johnson, editor, Proceedings of the First Annual ACM-SIAM Symposium on Discrete Algorithms, 22-24 January 1990, San Francisco, California, USA, pages 201–211. SIAM, 1990. URL: http://dl.acm.org/citation.cfm?id=320176.320199.

\bibitem[4]{bennett1993teleporting}
Charles H Bennett, Gilles Brassard, Claude Crépeau, Richard Jozsa, Asher Peres, and William K Wootters. Teleporting an unknown quantum state via dual classical and Einstein- Podolsky-Rosen channels. Physical review letters, 70(13):1895, 1993.

\bibitem[5]{BILPS:20}
Markus Bläser, Christian Ikenmeyer, Vladimir Lysikov, Anurag Pandey, and Frank-Olaf Schreyer. Variety membership testing, algebraic natural proofs, and geometric complexity theory. arXiv, 2020. arXiv:1911.02534.

\bibitem[6]{BL08}
Peter A. Brooksbank and Eugene M. Luks. Testing isomorphism of modules. Journal of Algebra, 320(11):4020–4029, 2008. Computational Algebra. doi:10.1016/j.jalgebra.2008.07.014.

\bibitem[7]{BCMW:17}
Peter Bürgisser, Matthias Christandl, Ketan D. Mulmuley, and Michael Walter. Membership in moment polytopes is in NP and coNP. SIAM J. Comput., 46(3):972–991, 2017. doi: 10.1137/15M1048859.

\bibitem[8]{BFGOWW2}
Peter Bürgisser, Cole Franks, Ankit Garg, Rafael Mendes de Oliveira, Michael Walter, and Avi Wigderson. Towards a theory of non-commutative optimization: geodesic first and second order methods for moment maps and polytopes. In 60th Annual IEEE Symposium on Foundations of Computer Science—FOCS 2019, pages 845–861. IEEE Computer Soc., Los Alamitos, CA, 2019. arXiv:1910.12375.

\bibitem[9]{BFGOWW}
Peter Bürgisser, Cole Franks, Ankit Garg, Rafael Oliveira, Michael Walter, and Avi Wigderson. Efficient algorithms for tensor scaling, quantum marginals, and moment polytopes. In 59th Annual IEEE Symposium on Foundations of Computer Science—FOCS 2018, pages 883–897. IEEE Computer Soc., Los Alamitos, CA, 2018. doi:10.1109/FOCS.2018.00088.

\bibitem[10]{BGOWW}
Peter Bürgisser, Ankit Garg, Rafael Oliveira, Michael Walter, and Avi Wigderson. Alternating minimization, scaling algorithms, and the null-cone problem from invariant theory. In 9th Innovations in Theoretical Computer Science, volume 94 of LIPIcs. Leibniz Int. Proc. Inform., pages Art. No. 24, 20. Schloss Dagstuhl. Leibniz-Zent. Inform., Wadern, 2018.

\bibitem[11]{BLNW:20}
Peter Bürgisser, Yinan Li, Harold Nieuwboer, and Michael Walter. Interior-point methods for unconstrained geometric programming and scaling problems. arXiv, 2020. arXiv:2008.12110.

\bibitem[12]{cohen2017matrix}
Michael B Cohen, Aleksander Madry, Dimitris Tsipras, and Adrian Vladu. Matrix scaling and balancing via box constrained Newton’s method and interior point methods. In Proceedings of the Symposium on Foundations of Computer Science (FOCS 2017), pages 902–913. IEEE, 2017. arXiv:1704.02310.

\bibitem[13]{cook:71}
Stephen A. Cook. The complexity of theorem proving procedures. In Proc. 3rd ACM STOC, pages 151–158, 1971.

\bibitem[14]{Cox-Little-Oshea}
David A. Cox, John Little, and Donal O’Shea. Ideals, varieties, and algorithms - an introduction to computational algebraic geometry and commutative algebra (2. ed.). Undergraduate texts in mathematics. Springer, 1997.

\bibitem[15]{Der01}
Harm Derksen. Polynomial bounds for rings of invariants. Proc. Amer. Math. Soc., 129(4):955– 963, 2001. doi:10.1090/S0002-9939-00-05698-7.

\bibitem[16]{Derksen-graph}
Harm Derksen. The graph isomorphism problem and approximate categories. J. Symb. Comput., 59:81–112, 2013. doi:10.1016/j.jsc.2013.06.002.

\bibitem[17]{DK}
Harm Derksen and Gregor Kemper. Computational invariant theory, volume 130 of Encyc- lopaedia of Mathematical Sciences. Springer, Heidelberg, enlarged edition, 2015. With two appendices by Vladimir L. Popov, and an addendum by Norbert A’Campo and Popov, Invariant Theory and Algebraic Transformation Groups, VIII. doi:10.1007/978-3-662-48422-7.

\bibitem[18]{DM-arbchar}
Harm Derksen and Visu Makam. Generating invariant rings of quivers in arbitrary character- istic. J. Algebra, 489:435–445, 2017. doi:10.1016/j.jalgebra.2017.06.035.

\bibitem[19]{DM-poly}
Harm Derksen and Visu Makam. Polynomial degree bounds for matrix semi-invariants. Adv. Math., 310:44–63, 2017. doi:10.1016/j.aim.2017.01.018.

\bibitem[20]{DM-si}
Harm Derksen and Visu Makam. Degree bounds for semi-invariant rings of quivers. J. Pure Appl. Algebra, 222(10):3282–3292, 2018. doi:10.1016/j.jpaa.2017.12.007.

\bibitem[21]{DM-oc}
Harm Derksen and Visu Makam. Algorithms for orbit closure separation for invariants and semi-invariants of matrices. Algebra Number Theory, 14(10):2791–2813, 2020. doi: 10.2140/ant.2020.14.2791.

\bibitem[22]{DM-exp}
Harm Derksen and Visu Makam. An exponential lower bound for the degrees of invariants of cubic forms and tensor actions. Adv. Math., 368:107136, 25, 2020. doi:10.1016/j.aim.2020. 107136.

\bibitem[23]{dolgachev:03}
Igor Dolgachev. Lectures on invariant theory, volume 296 of London Mathematical So- ciety Lecture Note Series. Cambridge University Press, Cambridge, 2003. doi:10.1017/ CBO9780511615436.

\bibitem[24]{durand-et-al:02}
Arnaud Durand, Miki Hermann, and Laurent Juban. On the complexity of recognizing the Hilbert basis of a linear Diophantine system. Theoret. Comput. Sci., 270(1-2):625–642, 2002. doi:10.1016/S0304-3975(01)00017-2.

\bibitem[25]{ESY14}
Kousha Etessami, Alistair Stewart, and Mihalis Yannakakis. A note on the complexity of comparing succinctly represented integers, with an application to maximum probability parsing. ACM Trans. Comput. Theory, 6(2):9:1–9:23, 2014. doi:10.1145/2601327.

\bibitem[26]{FS}
Michael A. Forbes and Amir Shpilka. Explicit Noether normalization for simultaneous conjug- ation via polynomial identity testing. In Approximation, randomization, and combinatorial optimization, volume 8096 of Lecture Notes in Comput. Sci., pages 527–542. Springer, Heidel- berg, 2013. doi:10.1007/978-3-642-40328-6\_37.

\bibitem[27]{GGOW16}
Ankit Garg, Leonid Gurvits, Rafael Oliveira, and Avi Wigderson. A deterministic polynomial time algorithm for non-commutative rational identity testing. In 57th Annual IEEE Symposium on Foundations of Computer Science—FOCS 2016, pages 109–117. IEEE Computer Soc., Los Alamitos, CA, 2016. doi:10.1109/FOCS.2016.95.

\bibitem[28]{GGOW:20}
Ankit Garg, Leonid Gurvits, Rafael Oliveira, and Avi Wigderson. Operator scaling: theory and applications. Found. Comput. Math., 20(2):223–290, 2020. doi:10.1007/ s10208-019-09417-z.

\bibitem[29]{GIMOWW}
Ankit Garg, Christian Ikenmeyer, Visu Makam, Rafael Mendes de Oliveira, Michael Walter, and Avi Wigderson. Search problems in algebraic complexity, GCT, and hardness of generators for invariant rings. In Shubhangi Saraf, editor, 35th Computational Complexity Conference, CCC 2020, July 28-31, 2020, Saarbrücken, Germany (Virtual Conference), volume 169 of LIPIcs, pages 12:1–12:17. Schloss Dagstuhl - Leibniz-Zentrum für Informatik, 2020. doi: 10.4230/LIPIcs.CCC.2020.12.

\bibitem[30]{Ge93}
Guoqiang Ge. Testing equalities of multiplicative representations in polynomial time (extended abstract). In 34th Annual Symposium on Foundations of Computer Science, Palo Alto, California, USA, 3-5 November 1993, pages 422–426. IEEE Computer Society, 1993. doi: 10.1109/SFCS.1993.366845.

\bibitem[31]{GLS-book}
Martin Grötschel, László Lovász, and Alexander Schrijver. Geometric algorithms and com- binatorial optimization, volume 2 of Algorithms and Combinatorics. Springer-Verlag, Berlin, second edition, 1993. doi:10.1007/978-3-642-78240-4.

\bibitem[32]{guillemin1990symplectic}
Victor Guillemin and Shlomo Sternberg. Symplectic techniques in physics. Cambridge university press, 1990.

\bibitem[33]{Gur04}
Leonid Gurvits. Classical complexity and quantum entanglement. J. Comput. Syst. Sci., 69(3):448–484, 2004. doi:10.1016/j.jcss.2004.06.003.

\bibitem[34]{gurvits2004combinatorial}
Leonid Gurvits. Combinatorial and algorithmic aspects of hyperbolic polynomials. arXiv preprint, 2004. arXiv:math/0404474.

\bibitem[35]{Hilb2}
D. Hilbert. Über die vollen Invariantensysteme. Math. Ann., 42:313–373, 1893. URL: http://eudml.org/doc/157652.

\bibitem[36]{Hilb1}
David Hilbert. Über die Theorie der algebraischen Formen. Math. Ann., 36(4):473–534, 1890. doi:10.1007/BF01208503.

\bibitem[37]{formal-ref-37}
Evelyne Hubert and George Labahn. Scaling invariants and symmetry reduction of dy- namical systems. Foundations of Computational Mathematics, 13, 2013. doi:10.1007/ s10208-013-9165-9.

\bibitem[38]{Humphreys:75}
James E. Humphreys. Linear algebraic groups. Springer-Verlag, New York-Heidelberg, 1975. Graduate Texts in Mathematics, No. 21.

\bibitem[39]{ikenmeyer2017vanishing}
Christian Ikenmeyer, Ketan D Mulmuley, and Michael Walter. On vanishing of Kronecker coefficients. Computational Complexity, 26(4):949–992, 2017.

\bibitem[40]{IQS}
Gábor Ivanyos, Youming Qiao, and K. V. Subrahmanyam. Non-commutative Edmonds’ problem and matrix semi-invariants. Comput. Complexity, 26(3):717–763, 2017. doi:10.1007/ s00037-016-0143-x.

\bibitem[41]{IQS2}
Gábor Ivanyos, Youming Qiao, and K. V. Subrahmanyam. Constructive non-commutative rank computation is in deterministic polynomial time. Comput. Complexity, 27(4):561–593, 2018. doi:10.1007/s00037-018-0165-7.

\bibitem[42]{KI04}
Valentine Kabanets and Russell Impagliazzo. Derandomizing polynomial identity tests means proving circuit lower bounds. Computational Complexity, 13(1-2):1–46, 2004. doi:10.1007/ s00037-004-0182-6.

\bibitem[43]{snf}
Ravindran Kannan and Achim Bachem. Polynomial algorithms for computing the Smith and Hermite normal forms of an integer matrix. SIAM J. Comput., 8(4):499–507, 1979. doi:10.1137/0208040.

\bibitem[44]{karp:72}
Richard M. Karp. Reducibility among combinatorial problems. In Complexity of computer computations (Proc. Sympos., IBM Thomas J. Watson Res. Center, Yorktown Heights, N.Y., 1972), pages 85–103, 1972.

\bibitem[45]{KN:78}
George Kempf and Linda Ness. The length of vectors in representation spaces. In Algebraic geometry (Proc. Summer Meeting, Univ. Copenhagen, Copenhagen, 1978), volume 732 of Lecture Notes in Math., pages 233–243. Springer, Berlin, 1979.

\bibitem[46]{kraft}
Hanspeter Kraft. Geometrische Methoden in der Invariantentheorie. Aspects of Mathematics, D1. Friedr. Vieweg \& Sohn, Braunschweig, 1984. doi:10.1007/978-3-322-83813-1.

\bibitem[47]{leep1999marriage}
David B Leep and Gerry Myerson. Marriage, magic, and solitaire. The American Mathematical Monthly, 106(5):419–429, 1999.

\bibitem[48]{levin:73}
L. A. Levin. Universal enumeration problems. Problemy Peredači Informacii, 9(3):115–116, 1973.

\bibitem[49]{linial2014vertices}
Nathan Linial and Zur Luria. On the vertices of the d-dimensional Birkhoff polytope. Discrete \& Computational Geometry, 51(1):161–170, 2014.

\bibitem[50]{linial2000deterministic}
Nathan Linial, Alex Samorodnitsky, and Avi Wigderson. A deterministic strongly polynomial algorithm for matrix scaling and approximate permanents. Combinatorica, 20(4):545–568, 2000.

\bibitem[51]{MW19}
Visu Makam and Avi Wigderson. Singular tuples of matrices is not a null cone (and, the symmetries of algebraic varieties). CoRR, abs/1909.00857, 2019. arXiv:1909.00857.

\bibitem[52]{GCTV}
Ketan D. Mulmuley. Geometric complexity theory V: Efficient algorithms for Noether normal- ization. J. Amer. Math. Soc., 30(1):225–309, 2017. doi:10.1090/jams/864.

\bibitem[53]{mulmuley2001geometric}
Ketan D Mulmuley and Milind Sohoni. Geometric complexity theory I: An approach to the P vs. NP and related problems. SIAM Journal on Computing, 31(2):496–526, 2001.

\bibitem[54]{mumford:88}
David Mumford. The red book of varieties and schemes, volume 1358 of Lecture Notes in Mathematics. Springer-Verlag, Berlin, 1988. doi:10.1007/978-3-662-21581-4.

\bibitem[55]{Mumford-book}
David Mumford, John Fogarty, and Frances Kirwan. Geometric invariant theory, Third Edition, volume 34 of Ergebnisse der Mathematik und ihrer Grenzgebiete. Springer, 1994.

\bibitem[56]{Popov-occ}
Vladimir L Popov. Two orbits: When is one in the closure of the other? Proceedings of the Steklov Institute of Mathematics, 264(1):146–158, 2009.

\bibitem[57]{schrijver:86}
Alexander Schrijver. Theory of linear and integer programming. Wiley-Interscience Series in Discrete Mathematics. John Wiley \& Sons, Ltd., Chichester, 1986. A Wiley-Interscience Publication.

\bibitem[58]{Sink}
R. Sinkhorn. A relationship between arbitrary positive matrices and doubly stochastic matrices. The Annals of Mathematical Statistics, 35:876–879, 1964.

\bibitem[59]{smith}
Henry J. Stephen Smith. On systems of linear indeterminate equations and congruences. Philosophical Transactions of the Royal Society of London, 151:293–326, 1861. URL: http: //www.jstor.org/stable/108738.

\bibitem[60]{Sturmfels}
Bernd Sturmfels. Algorithms in Invariant Theory. Texts \& Monographs in Symbolic Compu- tation. Springer, 2008. doi:10.1007/978-3-211-77417-5.

\bibitem[61]{Wehlau93}
David Wehlau. Constructive invariant theory for tori. Annales de l’institut Fourier, 43(4):1055– 1066, 1993. URL: http://eudml.org/doc/75025.

\bibitem[62]{weyl:39}
Hermann Weyl. The Classical Groups. Their Invariants and Representations. Princeton University Press, Princeton, N.J., 1939.
\end{thebibliography}
